\documentclass[a4paper,12pt]{article}

\newcommand*{\fg}[4][\textwidth]{
    \begin{figure}[!htb]
        \begin{center}
            \includegraphics[width=#1]{#2}
            \caption{#3}
            \label{rys:#4}
        \end{center}
    \end{figure}
}

\newcommand*{\Oznacz}[2]{\ref{#1:#2} (s. \pageref{#1:#2})}
\newcommand*{\OznaczZdjecie}[2][Rys.]{#1 \Oznacz{rys}{#2}}
\newcommand*{\OznaczKod}[1]{\Oznacz{lst}{#1}}
\newcommand*{\ListingFile}[3]{\lstinputlisting[caption=#1, label={lst:#2}, language=C++]{#3}}

\usepackage[utf8]{inputenc}
\usepackage[T1]{fontenc}
\usepackage{lmodern}
\usepackage{geometry}
\usepackage{titling}

\pretitle{\begin{center}\small{ZARZĄDZANIE PROJEKTEM}\\\vspace{1cm}\LARGE\bfseries}
\title{PLAN PROJEKTU}
\posttitle{\vspace{1cm}\\\large{\textit{PARKING PODZIEMNY}}\\[2cm]\end{center}}


\usepackage[polish]{babel}
\usepackage{amsmath}
\usepackage{amsfonts}
\usepackage{amssymb}
\usepackage{graphicx}
\usepackage{hyperref}
\usepackage{geometry}
\usepackage{caption}
\usepackage{listings}
\usepackage{xcolor}
\usepackage{fancyhdr}
\usepackage{hyperref}
\usepackage{multirow}
\usepackage{tabularx}

% New command for bottom dots
\newcommand{\bdots}{\vspace*{\fill}\dotfill}

\pagestyle{empty}
\geometry{margin=0.5in}
\date{}  % This will remove the date

\begin{document}
\maketitle

\begin{center}
\begin{tabularx}{\textwidth}{X X X X}
\hline
{\centering Imię i nazwisko\par} & {\centering EUU1\par} & {\centering EUU2\par} & {\centering EUK1\par} \\[5ex]
\hline
\rule{0pt}{5ex}
\multirow{2}{*}{Dominik Żuchowicz} & \bdots & \bdots & \bdots \\[1ex]
\textit{(kierownik projektu)} & & & \\[4ex]
Maciej Wójs & \bdots & \bdots & \bdots \\[4ex]
Jakub Tokarczyk & \bdots & \bdots & \bdots \\[4ex]
Arkadiusz Ryczek & \bdots & \bdots & \bdots \\[4ex]
\end{tabularx}
\end{center}
\texttt{Informatyka, I stopień, 3 rok\\
Grupa nr: 4\\
Studia: stacjonarne\\
}
\clearpage
\tableofcontents
\clearpage
\section{Uzasadnienie projektu}
Obecna dostępność miejsc parkingowych wokół budynku uczelni jest niewystarczająca — powoduje opóźnienia, frustrację użytkowników, parkowanie w niedozwolonych miejscach i zwiększone natężenie ruchu na drogach dojazdowych. Realizacja parkingu ma poprawić ergonomię dojazdu, zmniejszyć chaos komunikacyjny i wesprzeć rozwój infrastruktury uczelni.

\section{Cele projektu (SMART)}
\begin{itemize}
    \item \textbf{S (Specific)} – Zbudowanie kompleksu parkingowego pod lub obok budynku uczelni w celu zwiększenia liczby miejsc postojowych.
    \item \textbf{M (Measurable)} – Utworzenie co najmniej 150 nowych miejsc parkingowych oraz skrócenie średniego czasu poszukiwania miejsca o 40\%.
    \item \textbf{A (Achievable)} – Realizacja projektu w oparciu o dostępny teren uczelni oraz możliwe źródła finansowania (dotacje, partner prywatny).
    \item \textbf{R (Relevant)} – Projekt odpowiada na realne potrzeby studentów i pracowników oraz wspiera wizerunek uczelni jako instytucji dbającej o infrastrukturę.
    \item \textbf{T (Time-bound)} – Zakończenie inwestycji w ciągu 18 miesięcy od rozpoczęcia prac budowlanych.
\end{itemize}

\section{Zakres projektu}
Projekt obejmuje:
\begin{itemize}
    \item Analizy wstępne: geotechniczne, ruchowe i środowiskowe.
    \item Wybór wariantu budowy (parking podziemny lub naziemny).
    \item Opracowanie dokumentacji projektowej (budowlanej i wykonawczej).
    \item Uzyskanie pozwoleń administracyjnych.
    \item Realizację robót budowlanych wraz z instalacjami.
    \item Montaż systemu zarządzania parkingiem (kontrola dostępu, monitoring, system płatności).
    \item Testy, odbiory techniczne i przekazanie obiektu do użytkowania.
\end{itemize}

\section{Interesariusze}
\begin{itemize}
    \item \textbf{Zarząd uczelni / Rektor} – decydent strategiczny, zatwierdza budżet i harmonogram.
    \item \textbf{Kierownik projektu (PM)} – odpowiedzialny za całościowe prowadzenie projektu.
    \item \textbf{Architekt (ARCH)} – przygotowuje koncepcję architektoniczną.
    \item \textbf{Inżynier konstrukcji (STR)} – opracowuje projekt konstrukcyjny.
    \item \textbf{Geotechnik (GEO)} – wykonuje badania gruntu.
    \item \textbf{Specjalista ds. środowiska (ENV)} – sporządza ocenę oddziaływania na środowisko.
    \item \textbf{Specjalista ds. ruchu (TRAF)} – analizuje natężenie ruchu i drogi dojazdowe.
    \item \textbf{Prawnik (PRAW)} – prowadzi sprawy formalno-prawne i pozwolenia.
    \item \textbf{Kosztorysant (KOSZT)} – opracowuje budżet i analizę ekonomiczną.
    \item \textbf{Wykonawca (WYK)} – realizuje prace budowlane.
    \item \textbf{Specjalista IT (IT)} – wdraża systemy parkingowe.
    \item \textbf{Inspektor BHP (BHP)} – nadzoruje bezpieczeństwo prac.
    \item \textbf{Społeczność lokalna / studenci / pracownicy} – użytkownicy końcowi parkingu.
\end{itemize}

\section{WBS -- Struktura podziału pracy}
\begin{enumerate}
    \item Inicjacja projektu
    \begin{itemize}
        \item Powołanie zespołu projektowego.
        \item Studium wykonalności i analiza wariantów.
        \item Zatwierdzenie budżetu i harmonogramu.
    \end{itemize}
    \item Planowanie
    \begin{itemize}
        \item Opracowanie szczegółowego harmonogramu (MS Project).
        \item Przygotowanie planu ryzyka i komunikacji.
        \item Ustalenie źródeł finansowania.
    \end{itemize}
    \item Projektowanie i uzgodnienia
    \begin{itemize}
        \item Badania geotechniczne i środowiskowe.
        \item Projekt architektoniczny i konstrukcyjny.
        \item Projekt systemu zarządzania parkingiem.
        \item Uzyskanie pozwoleń administracyjnych.
    \end{itemize}
    \item Realizacja
    \begin{itemize}
        \item Roboty ziemne i konstrukcyjne.
        \item Instalacje elektryczne, sanitarne i IT.
        \item Wykończenie, oznakowanie, zieleń.
    \end{itemize}
    \item Testy i odbiory
    \begin{itemize}
        \item Odbiory techniczne i BHP.
        \item Testy systemów parkingowych.
        \item Szkolenie operatorów i przekazanie obiektu.
    \end{itemize}
\end{enumerate}

\section{Harmonogram projektu i zasoby}
\begin{itemize}
    \item Całkowity czas trwania projektu: 18–24 miesiące.
    \item Główne zasoby:
    \begin{itemize}
        \item Kadrowe: PM, ARCH, STR, GEO, ENV, TRAF, IT, WYK, BHP.
        \item Materialne: sprzęt budowlany, materiały konstrukcyjne, systemy parkingowe (bramki, czujniki, CCTV).
        \item Finansowe: budżet inwestycyjny (środki uczelni, dotacje, PPP).
    \end{itemize}
\end{itemize}

\section{Budżet projektu}
\begin{itemize}
    \item Szacunkowy koszt: 8–12 mln PLN (zależnie od wariantu konstrukcji).
    \item Rezerwa budżetowa: 10\% wartości inwestycji.
    \item Źródła finansowania:
    \begin{itemize}
        \item Środki własne uczelni.
        \item Dotacje miejskie / unijne.
        \item Partner prywatny (model PPP).
    \end{itemize}
\end{itemize}

\section{Komunikacja w zespole projektowym}
\begin{itemize}
    \item Cotygodniowe spotkania statusowe zespołu projektowego.
    \item Miesięczne raporty postępu do władz uczelni.
    \item Kanały komunikacji: Teams / e-mail / tablica MS Project.
    \item Rejestr decyzji projektowych (Project Decision Log).
    \item Konsultacje społeczne z mieszkańcami i użytkownikami kampusu.
\end{itemize}

\section{Zarządzanie ryzykiem}
\subsection*{Macierz ryzyka}
\begin{center}
\renewcommand{\arraystretch}{1.3}
\begin{tabular}{|p{4.5cm}|c|c|p{5cm}|}
\hline
\textbf{Ryzyko} & \textbf{Prawdopodobieństwo} & \textbf{Skutek} & \textbf{Działania zaradcze} \\
\hline
Nieodpowiednie warunki geotechniczne (np. wysoki poziom wód gruntowych) & Wysokie & Wysokie & Dodatkowe badania gruntu, możliwość zmiany wariantu konstrukcji na naziemny. \\
\hline
Opóźnienia administracyjne w uzyskaniu pozwoleń & Średnie & Wysokie & Wczesne rozpoczęcie procedur, ścisła współpraca z urzędem miasta. \\
\hline
Przekroczenie budżetu inwestycji & Średnie & Wysokie & Stały nadzór kosztorysanta, rezerwa finansowa 10\%. \\
\hline
Protesty mieszkańców / społeczności lokalnej & Niskie & Średnie & Konsultacje społeczne, komunikacja o korzyściach projektu. \\
\hline
Awaria systemu parkingowego po uruchomieniu & Niskie & Wysokie & Testy integracyjne przed odbiorem, umowa serwisowa z dostawcą IT. \\
\hline
Utrudnienia w ruchu podczas budowy & Średnie & Średnie & Opracowanie tymczasowej organizacji ruchu, informowanie użytkowników uczelni. \\
\hline
Brak dostępnych wykonawców w przetargu & Niskie & Wysokie & Analiza rynku wcześniej, elastyczne terminy przetargu. \\
\hline
\end{tabular}
\end{center}

\section{Analiza SWOT}
\subsection*{Mocne strony (Strengths)}
\begin{itemize}
    \item Rozwiązanie realnego problemu braku miejsc parkingowych.
    \item Wsparcie władz uczelni i zgodność z jej strategią rozwoju.
    \item Możliwość zastosowania nowoczesnych technologii (systemy IoT, płatności automatyczne).
    \item Poprawa bezpieczeństwa i przepustowości komunikacyjnej.
\end{itemize}

\subsection*{Słabe strony (Weaknesses)}
\begin{itemize}
    \item Wysokie koszty inwestycyjne.
    \item Ryzyko techniczne i geotechniczne (szczególnie dla wariantu podziemnego).
    \item Długi czas uzyskiwania pozwoleń.
    \item Potencjalne utrudnienia w ruchu podczas budowy.
\end{itemize}

\subsection*{Szanse (Opportunities)}
\begin{itemize}
    \item Możliwość pozyskania finansowania z funduszy miejskich lub unijnych.
    \item Partnerstwo publiczno-prywatne (PPP).
    \item Wzrost wizerunku uczelni jako nowoczesnej instytucji.
    \item Integracja z miejskimi systemami mobilności.
\end{itemize}

\subsection*{Zagrożenia (Threats)}
\begin{itemize}
    \item Opóźnienia formalne lub prawne.
    \item Wzrost kosztów materiałów budowlanych.
    \item Problemy techniczne przy realizacji.
    \item Niska akceptacja społeczna dla budowy w sąsiedztwie kampusu.
\end{itemize}

\end{document}


